\section{Detailed Experimental Settings}~\label{appx:settings}

This appendix specifies the planner settings, query policies, domain models, and reporting conventions for Sec.~\ref{exp:setup}. The oracle experiments use simulated execution with correct responses. The VLM and human evaluations use physical robot execution. We aggregate trials within each domain for all reported comparisons.

\subsection{Planner and Evaluation Parameters}~\label{appx:planner_params}
Table~\ref{tab:app_planner_params} lists the default POMCP parameters for Ours. The oracle trial limit is $50$ physical actions, and the physical trial limit is $20$. Queries do not consume physical steps for Ours. Query-Action uses a limit of $50$ total decisions that includes queries and physical actions. \cp{Seven WasteSorting trials and $13$ TomatoHarvesting trials terminate at this limit.} The reward-based configurations use the discount factors specified below.

\begin{table}[!htbp]
\centering
\caption{Planner and evaluation parameters used by Ours.}
\label{tab:app_planner_params}
\small
\begin{tabular}{lc}
\toprule
\textbf{Parameter} & \textbf{Value} \\
\midrule
POMCP simulations per planning step & $100$ \\
Maximum simulation depth & $20$ \\
Discount factor $\gamma$ & $0.2$ \\
UCB exploration constant $c$ & $1.0$ \\
Oracle physical action limit & $50$ \\
Physical action limit & $20$ \\
Default confidence threshold $\tau$ & $0.8$ \\
Discount cutoff $\epsilon$ & $0.005$ \\
Oracle frontier capacity & $8{,}000$ \\
Oracle search-node particle capacity & $8{,}000$ \\
Oracle transition sampling budget & $250$ \\
Oracle initial configurations per domain & $5$ \\
Trials per initial configuration and condition & $40$ \\
Oracle trials per domain and condition & $200$ \\
\bottomrule
\end{tabular}
\end{table}

\cp{With $\gamma=0.2$ and $\epsilon=0.005$, the discount cutoff stops search at depth $4$ before the nominal maximum of $20$.}

\subsection{Baseline Configurations}~\label{appx:baseline_configs}
The oracle baseline comparison evaluates KnowNo, IntroPlan, and Query-Action against Ours. The physical experiments compare KnowNo and Ours. Appendix~\ref{appx:baseline_algorithms} provides separate procedures for the three baselines. KnowNo and IntroPlan use GPT-4o for action generation and candidate scoring. This planning model is separate from Gemma 4 31B IT, which provides external feedback in the VLM experiments.

The \cp{recorded} action generation temperature is $0.3$, and the candidate score temperature is $5.0$. KnowNo uses calibration thresholds $\hat{q}=0.845165$ for WasteSorting and $\hat{q}=\cp{0.8404}$ for TomatoHarvesting in the oracle comparison. The physical VLM logs \cp{use $0.845165$ and $0.901119$ respectively}, except for one WasteSorting trial with $\hat{q}=0.8704$. \cp{The oracle TomatoHarvesting KnowNo batch does not record the generation temperature.} IntroPlan retrieves three examples and uses calibration thresholds $0.980493$ and $0.985552$ for the two domains in the same order. Query-Action uses the \cp{recorded domain-specific configuration}. The discount factor is $0.9$ for WasteSorting and $0.5$ for TomatoHarvesting. Query cost is zero, the failure penalty is $10$, assumed response accuracy is $1.0$, and the planner uses $100$ simulations per decision.

% Source: experiments_logs/query_action_tuning/20260921_130229_529697/analysis/
% report.md, selected_config.json, and 01_processed/episodes.csv (1,200 trials).
A preliminary Query-Action sweep tested $\gamma\in\{0.2,0.5,0.8,0.95\}$ and query costs $\{0,0.5,1\}$, with ten trials per scene in each domain for every pair. The selection criterion was the highest pooled task success rate: $\gamma=0.5$ with zero query cost achieved $66/100$ successes. This supports the zero-cost setting and the TomatoHarvesting discount factor used in the reported configurations.
% TODO(author): Explain the later WasteSorting choice gamma=0.9. It is verified
% in the final execution logs but is absent from the above sweep; the sweep
% winner was gamma=0.5 for the pooled objective, not a domain-specific optimum.

\subsection{Ablation Configurations}~\label{appx:ablation_configs}
We distinguish the criterion for entering a query episode, the rule for selecting each fact, and the criterion for ending the episode. These decisions follow the belief update from a physical observation and precede the next physical action.

\noindent\textbf{Entry.} Fully-Random and W1 enter with probability $0.4$ per physical step. Ours, Q1, and Q2 enter when belief confidence is below $\tau=0.8$. W2 enters unless the conformal action prediction set contains exactly one concrete action; an empty set or a set containing the fallback option also triggers verification. W3 enters when the highest query-action value exceeds the highest applicable physical-action value.

\noindent\textbf{Selection.} Fully-Random and Q1 sample uniformly from facts that distinguish the remaining state hypotheses. Ours and W1--W3 select the fact with the largest expected entropy reduction, recomputed after each response. Thus neither W2's conformal set nor W3's highest-valued query specifies the fact actually asked. Q2 instead selects the highest-valued query fact and recomputes query values after each response; it does not compare them with physical-action values to cancel a query.

% The 396 valid archived W2 policy traces, not the current query_episode.py,
% verify repeated CP decisions. W3 uses the 20261005_204332_446690 rerun;
% Q2 uses the 20261005_190346_519071 batch. See asset/BRL_oracle_exp_v2/
% 02_when_what_policy_ablation/00_raw/ and EXPERIMENT_SETTINGS.md.
\noindent\textbf{Continuation and termination.} After entry, Fully-Random, W1, W3, Q1, Q2, and Ours continue until confidence reaches $\tau$ or no differentiating fact remains. W3 evaluates reward-based timing only on entry and has no one-question limit in the reported rerun. In contrast, the archived W2 batch used for the reported results reevaluates the conformal criterion after each fact response and stops when the set becomes a singleton concrete action or no differentiating fact remains. W2 can therefore stop below $\tau$; its continuation criterion differs from confidence-based verification. A failure to reduce the represented hypotheses also ends an episode in the policy runner.

W2 uses calibration thresholds $0.8704$ for WasteSorting and $0.8404$ for TomatoHarvesting, with score temperature $5.0$. Four unparseable outputs count as failures in the success rate and contribute no successful-trial cost measurements. Both reward-based variants use zero query cost, a failure penalty of $10$, and assumed response accuracy $1.0$. They reuse the domain-specific discount factors of Query-Action, $0.9$ for WasteSorting and $0.5$ for TomatoHarvesting, for both query evaluation and physical action planning. The query evaluator uses at least $100$ simulations and increases the budget to the number of root candidates plus one if necessary. The remaining configurations use $\gamma=0.2$. Consequently, comparisons with W3 and Q2 also change the discount factor, and the W2 comparison changes the continuation rule in addition to the entry criterion.

\subsection{Oracle Reference Trials}~\label{appx:oracle_reference}
We use the complete Ours batch from the policy ablation as a shared oracle reference for the baseline and ablation comparisons as well as the threshold analysis at $\tau=0.8$. This reference contains $200$ trials per domain and is reused across these comparisons. The other threshold points retain their original trial batches.

% Parameter values below are verified against the current 02 domain YAML files.
% Historical per-run YAML snapshots are not included in the raw bundles.
\subsection{Domain Model Probabilities}~\label{appx:model_params}
Table~\ref{tab:app_model_params} lists the parameters in the domain configuration files. Transition probabilities determine symbolic successor hypotheses. Observation sampling generates simulated sensor responses, and observation likelihoods weight the represented hypotheses. Detection sampling accuracy and likelihood accuracy are specified separately. Physical execution supplies real sensor observations to the belief update.

\begin{table}[!htbp]
\centering
\caption{Domain model parameters for transition and observation models.}
\label{tab:app_model_params}
\small
\begin{tabular}{lll}
\toprule
\textbf{Domain} & \textbf{Model component} & \textbf{Probability} \\
\midrule
WasteSorting & Detect observed transition & $0.90$ \\
WasteSorting & Detect category transition & $0.50$ \\
WasteSorting & Pick transition & $0.90$ \\
WasteSorting & Place transition & $0.90$ \\
WasteSorting & Detect observation likelihood & $0.97$ \\
WasteSorting & Detect category observation & $0.85$ \\
WasteSorting & Pick/place observation & $0.95$ \\
WasteSorting & Detect sampling accuracy & $[0.94,1.00]$ \\
WasteSorting & False-positive likelihood & $0.03$ \\
\midrule
TomatoHarvesting & Navigate transition & $0.99$ \\
TomatoHarvesting & Detect transition & $0.95$ \\
TomatoHarvesting & Pick transition & $0.95$ \\
TomatoHarvesting & Scan transition & $0.90$ \\
TomatoHarvesting & Place/discard transition & $0.99$ \\
TomatoHarvesting & Detect observation likelihood & $0.90$ \\
TomatoHarvesting & Detect classification observation & $0.95$ \\
TomatoHarvesting & Scan observation & $0.85$ \\
TomatoHarvesting & Detect sampling accuracy & $[0.85,0.95]$ \\
TomatoHarvesting & False-positive probability & $0.03$ \\
TomatoHarvesting & Navigate observation & $0.95$ \\
TomatoHarvesting & Place observation & $1.00$ \\
\bottomrule
\end{tabular}
\end{table}

The WasteSorting category transition parameter $0.50$ creates competing category hypotheses after detection. Detection observation sampling uses the hidden object category, with uniform allocation of classification error across the other categories. Occluded objects and objects already held or placed in bins are excluded from simulated detection. In TomatoHarvesting, detection produces ripe or unripe labels, and rotten tomatoes share the ripe visual label. A subsequent scan distinguishes fresh from rotten tomatoes. The detect models combine individual object outcomes into a joint observation distribution. Detector confidence values do not directly scale the symbolic observation likelihood.

\subsection{Randomization and Episode Seeds}~\label{appx:randomization}
Each oracle trial records a pseudorandom seed. The runner uses the seed for Python \texttt{random}, NumPy sampling, initial arrangements, and randomized query decisions. Conditions within each oracle experiment share the seed assigned to the same domain, initial configuration, and repetition. Different policies can consume random numbers differently after different action choices. The shared Ours reference described in Appendix~\ref{appx:oracle_reference} retains its original seeds.

\subsection{Uncertainty Estimates and Time Measures}~\label{appx:metric_details}
Success rates use all attempted trials. Query counts, query probabilities, plan lengths, and time measures use successful trials and are reported as mean $\pm$ sample standard deviation. Query probability is the fraction of physical steps that initiate at least one query. Several questions within one query episode increase query count but contribute only one triggered step. Success bars show $95\%$ Wilson confidence intervals. Cost error bars show sample standard deviations. A dash or the asterisk identified in the figure caption denotes a single successful trial for which sample standard deviation is unavailable.

Search time is the sum of directly recorded planning durations within a trial. The oracle plan-time comparison uses cumulative POMCP search time for Ours and Query-Action and cumulative candidate generation and scoring time for KnowNo and IntroPlan. Interaction time measures external requests and response processing and is reported for the physical experiments. Total time is the accumulated trial duration. Physical result tables and figures use the residual planning-time estimate
\[
t_{\mathrm{plan}}=t_{\mathrm{total}}-t_{\mathrm{interaction}}-t_{\mathrm{execution}}.
\]
This estimate includes overhead outside the separately recorded interaction and physical execution durations. Direct search time and residual planning time are distinct measurements.

\newtcolorbox{vlmpromptbox}[1]{
    enhanced,
    breakable,
    sharp corners,
    colback=gray!12,
    colframe=black,
    colbacktitle=black,
    coltitle=white,
    fonttitle=\bfseries,
    fontupper=\small,
    boxrule=0.4pt,
    left=3mm,
    right=3mm,
    top=2mm,
    bottom=2mm,
    before skip=10pt,
    after skip=10pt,
    title={#1},
    title after break={}
}

\subsection{VLM Prompts}~\label{appx:vlm_prompts}

Both domains use the same external VLM configuration. The served model identifier is \texttt{nvidia/Gemma-4-31B-IT-NVFP4}. The sampling temperature is $0.3$, the maximum completion length is $256$ tokens, and the request timeout is $90$ seconds. Log probabilities are enabled with the top five token alternatives. Thinking and reasoning output are disabled. A text-only warmup asks whether the model is ready and requests a Yes or No response.

Each experimental request includes two current RGB camera images and the current robot action. Ours supplies a Korean state-verification question, the judgment criterion for the queried predicate, and the domain action model. KnowNo supplies candidate actions with response tokens, judgment criteria for the candidate action types, and the domain action model. TomatoHarvesting additionally supplies an environment model with fixed object identifiers and stem assignments. The following English descriptions summarize the configured prompt instructions. The supplementary configuration files preserve the original English and Korean prompt text.

\begin{vlmpromptbox}{Shared Request Instructions}
For state verification, the VLM answers the supplied Korean question from the two camera images and robot context. The response contains two lines, with Yes or No first and a concise reason second. The current action describes the attempted operation and does not establish that the queried fact is true. The VLM must judge the fact from visible evidence. The client converts Yes to \texttt{true} and No to \texttt{false} before the belief update.

For action selection, the VLM chooses the most appropriate next action from the supplied candidates. The response contains two lines, with a candidate token first and a concise Korean reason second. The VLM must verify the relevant object, robot state, and execution conditions from the images. The candidate name and current action do not establish feasibility. If no concrete candidate is appropriate, the VLM selects the option that requests a manually entered action. The client parses the selected token and checks that the token belongs to the supplied candidate set.

\end{vlmpromptbox}

\begin{vlmpromptbox}{WasteSorting Domain Instructions}
The action model specifies \texttt{detect}, \texttt{pick}, \texttt{place}, and \texttt{done}. Detection requires an empty gripper and identifies visible waste objects and categories. Picking requires a detected object and an empty gripper. Placement requires the robot to hold the selected object and use the bin corresponding to the detected category. Completion requires all waste objects to be sorted and the gripper to be empty. A visible object, bounding box, or category label does not establish that detection is complete. An object exposed after another object is removed must be detected before picking.

State-verification criteria use the identity and appearance of the specified object. For \texttt{detected}, the VLM first identifies the specified waste number and checks whether the object is visible on the worktable. Another visible object or the current detection action is insufficient evidence. Occlusion or uncertain identity must not be treated as proof that the object is absent. For category predicates, the VLM checks material and shape. Plastic requires visible features of a plastic bottle or container. A can requires the configured rectangular shape and metal surface, and a round plastic container must not be classified as a can. Paper requires features such as a thin surface, folds, or printed cardboard. The configured general-waste example has a banana shape. Color alone or the bin name in an action candidate is insufficient evidence of category.

For \texttt{holding} and \texttt{handempty}, the VLM checks the actual contents of the gripper in both images. For \texttt{in\_bin}, the specified object must be visibly inside the specified bin. For \texttt{sorted}, the object category and bin category must match. For \texttt{done}, no waste may remain on the worktable, all waste must be in the correct bins, and the gripper must be empty. Waste, robot, and bin-type predicates require visual identification of the specified target. Attempted pick or place actions do not establish the corresponding state facts.

The action-selection criteria require \texttt{detect} before \texttt{pick} for newly exposed or undetected waste. A pick candidate requires an empty gripper and a visible object with an identifiable number. A place candidate requires the gripper to hold the specified object and the visual category to match the target bin. A completion candidate requires an empty worktable, correct bin placement, and an empty gripper. If visibility, object identity, or execution conditions do not support a concrete candidate, the VLM selects the option for manual action entry.

\end{vlmpromptbox}

\begin{vlmpromptbox}{TomatoHarvesting Domain Instructions}
The environment model specifies two tomatoes per stem. Tomatoes 1 and 2 belong to \texttt{stem\_01}, and tomatoes 3 and 4 belong to \texttt{stem\_02}. Numbers run from right to left across the scene and remain fixed after occlusion, picking, harvesting, or disposal. The right pair of ArUco markers identifies \texttt{stem\_01}, and the left pair identifies \texttt{stem\_02}. The location criterion also states that a view containing only two markers indicates that the robot is at \texttt{stem\_02}.

The action model specifies navigation, detection, picking, scanning, placement, disposal, and completion. Navigation and detection require an empty gripper. The robot moves to the next stem if no harvestable tomato remains at the current stem. Picking requires an observed ripe tomato at the current stem and an empty gripper. Scanning inspects the held tomato. Placement requires a tomato confirmed fresh after scanning, and disposal requires a tomato confirmed rotten after scanning. Completion requires all necessary tomatoes to be processed. A newly exposed tomato requires detection before picking even if the tomato is visible in the images.

For \texttt{at}, the VLM identifies the specified tomato and target stem using the fixed identifiers and markers. Direct visibility at the target stem supports Yes, and direct visibility at the opposite stem supports No. Absence from an image or uncertain identity does not establish that the tomato belongs to another stem. A tomato held by the robot is no longer on the stem, and each stem contains at most two tomatoes. The current action must not be used as location evidence. For \texttt{ripe} and \texttt{unripe}, the dominant red or green surface color determines ripeness. Scratches and rotten areas are separate from ripeness.

For \texttt{fresh} and \texttt{rotten}, inspect the surface of the specified tomato actually held in the gripper. Answer No only if a large, distinct beige rotten area is visible. Small scratches, dents, spots, reflections, and minor discoloration near the stalk do not indicate rot and must not be used to reject freshness. Base the judgment on both the scene image and the wrist-camera image.

For \texttt{holding}, \texttt{holded}, and \texttt{handempty}, the VLM checks the actual gripper contents. For \texttt{loaded} and \texttt{discarded}, the specified tomato must be visibly in the harvest basket or disposal location. The predicates \texttt{observed} and \texttt{scanned} describe completed procedures. Visible tomatoes and attempted detect or scan actions do not establish procedural completion. For \texttt{located}, the VLM checks the robot position relative to the target marker region. Tomato counts use distinct visible tomatoes across both images without double counting or assumptions about occluded regions. Harvested counts use tomatoes visibly inside the basket and do not substitute the number of attempted placements. Tomato, robot, location, and stem predicates require identification from appearance, scene structure, and markers.

The action-selection criteria require the VLM to verify robot location, gripper contents, and the target tomato. Navigation is appropriate if the gripper is empty and movement to another stem is necessary. Detection precedes picking for newly visible or undetected tomatoes at the current stem. Picking requires an empty gripper and a visible ripe tomato at the current stem. Scanning requires the robot to hold the specified tomato and still need a quality judgment. Placement requires the held tomato to appear fresh without a large rotten area. Disposal requires a large, distinct rotten area and must not follow from minor scratches alone. Completion requires no remaining ripe tomatoes to process and an empty gripper. If no concrete action satisfies the visible conditions, the VLM selects the option for manual action entry.
\end{vlmpromptbox}

